\section{Домейн модел и модел на данните}

Домейн моделът е централната абстракция, чрез която софтуерната система представя предметната област. В разглеждания проект той описва основните участници и артефакти в процеса на електронна търговия: потребител, фирма, продукт, количка и история на поръчки. Макар приложението да не реализира сложна бизнес семантика от enterprise клас, домейн моделът е достатъчно богат, за да наложи ясно картографиране към релационна база данни и разделяне на различни типове модели за обработка и визуализация на данни.

Настоящата глава разглежда едновременно логическото съдържание на домейн обектите и начина, по който те са реализирани чрез JPA entity класове, binding модели, service модели и view модели. По този начин се проследява целият път на информацията: от въвеждането ѝ във форма, през вътрешното ѝ представяне в системата, до записа ѝ в базата и повторната ѝ визуализация.

\subsection{Основни домейн понятия}

В предметната област на електронния магазин могат да бъдат идентифицирани няколко базови понятия.

\begin{itemize}[leftmargin=*]
    \item \textbf{Потребител} -- физическо лице, което се регистрира, влиза в системата, поддържа профил, добавя фирми, въвежда продукти, използва количка и има история на поръчки.
    \item \textbf{Фирма} -- обект, който представя търговска или доставна организация, асоциирана с конкретен потребител и използвана при добавяне на продукти.
    \item \textbf{Продукт} -- артикул от каталога, описан чрез име, цена, наличност, изображение и категория.
    \item \textbf{Запис в количка} -- конкретно желание за покупка на определен продукт в дадено количество.
    \item \textbf{Исторически запис} -- обобщена информация за приключена поръчка, включваща адрес, дата и цена.
\end{itemize}

Тези понятия образуват ядрото на домейна и определят по-голямата част от класовете в \texttt{models.entities}. От гледна точка на \textit{domain-driven} мисленето \cite{evans}, те могат да бъдат разглеждани като основните бизнес обекти, около които се организират и останалите механизми в приложението.

\subsection{Entity моделът в JPA слоя}

Persistence слоят използва JPA entity класове, които наследяват \texttt{BaseEntity}. Макар в текущата работа да не е необходимо да се навлиза в детайли на базовия клас, неговото наличие подсказва стандартизирано централизиране на идентификатора за всички трайни обекти. Това е типична практика и намалява дублирането на общи persistence характеристики.

\subsubsection{Потребител}

Класът \texttt{User} е основният aggregate обект в системата. Той съдържа име, фамилия, парола, признак дали акаунтът е потвърден и електронна поща. Освен това поддържа връзки към фирмите на потребителя, към записите в количката и към историята на поръчките. Така именно потребителят е логическият център на множество операции в приложението.

Особен интерес представлява полето за електронна поща. То е маркирано като уникално и допълнително има регулярен израз на ниво entity, което подсказва стремеж към защита на едно от най-важните идентифициращи полета. Паролата не се валидира на това ниво, а се обработва и хешира чрез service логика. Булевото поле за потвърждение е част от модела на регистрация и контролира допустимостта на входа.

\subsubsection{Фирма}

Класът \texttt{Company} описва фирмена информация и включва атрибути като име, \texttt{mol}, \texttt{vat}, адрес и признак за регистрация. В системата фирмата принадлежи на конкретен потребител чрез връзка \texttt{@ManyToOne}. Освен това една фирма има множество продукти чрез \texttt{@OneToMany(mappedBy = "company")}. От логическа гледна точка това е разумен модел: фирмата изпълнява ролята на притежател или източник на предлаганите продукти.

Следва да се отбележи, че името на колекцията в класа \texttt{Company} е \texttt{companies}, въпреки че тя реално съдържа продукти. Това е незначително от функционална гледна точка, но е показателно за важността на консистентното именуване в домейн модела. Подобни дребни несъответствия не чупят системата, но намаляват концептуалната яснота и следва да бъдат разглеждани като потенциален обект на бъдещо рефакториране.

\subsubsection{Продукт}

Класът \texttt{Product} е основният каталожен обект в приложението. Той съдържа име, цена, количество, URL към изображение, категория и връзка към фирма. Категорията е реализирана като изброим тип \texttt{CategoryEnum}, което е добър архитектурен избор. По този начин наборът от допустими категории е изрично описан в кода и не зависи от свободен текст или произволни стойности, въведени от потребителя.

Валидирането на цена и количество е реализирано чрез \texttt{@Min(0)}, което формализира базовото бизнес правило, че продуктите не могат да имат отрицателни стойности за тези полета. От бизнес гледна точка продуктът е едновременно обект на визуализация и на транзакционно взаимодействие, защото участва в количката и се променя при checkout процеса.

\subsubsection{Запис в количка}

Класът \texttt{ShoppingCartEntity} моделира преходно търговско състояние: конкретен потребител е избрал конкретен продукт в конкретно количество. Той съдържа само три основни атрибута -- желано количество, продукт и потребител. Макар моделът да е опростен, той е достатъчен за целите на проекта и илюстрира идеята, че количката не е просто списък от продукти, а списък от връзки „потребител--продукт--брой“.

Интересен детайл е, че продуктът в количката е реализиран с \texttt{@OneToOne}, а потребителят с \texttt{@ManyToOne}. Това позволява множество записи в количката да принадлежат на един потребител, но ограничава логическата връзка към конкретен продукт. В production сценарий често би се използвала и различна схема за моделиране на поръчковите позиции, но за настоящия проект този модел е достатъчен, за да представи основната идея.

\subsubsection{История на поръчките}

Класът \texttt{HistoryEntity} представя обобщен запис за реализирана поръчка. В него се съхраняват адрес, дата, цена и връзка към потребител. Това е умишлено опростен модел. Вместо поръчката да се разглежда като самостоятелен aggregate с отделни позиции, статуси и платежни детайли, тук тя е сведена до исторически запис за случило се checkout събитие. Подобен подход е напълно оправдан в учебно-приложен контекст, тъй като позволява системата да покаже завършен цикъл на покупка без да навлиза в сложността на пълномащабен order management модел.

\begin{figure}[H]
    \centering
    \begin{tikzpicture}
        \node[appblock] (user) {User};
        \node[appblock, below left=2.2cm and 2.3cm of user] (company) {Company};
        \node[appblock, below right=2.2cm and 2.3cm of user] (history) {HistoryEntity};
        \node[appblock, below=2.2cm of user] (cart) {ShoppingCartEntity};
        \node[appblock, below=2.2cm of company] (product) {Product};

        \draw[appline] (company) -- node[left]{Many-to-One} (user);
        \draw[appline] (history) -- node[right]{Many-to-One} (user);
        \draw[appline] (cart) -- node[right]{Many-to-One} (user);
        \draw[appline] (product) -- node[left]{Many-to-One} (company);
        \draw[appline] (cart) -- node[below]{One-to-One} (product);
    \end{tikzpicture}
    \caption{Обобщена схема на връзките между основните entity обекти}
    \label{fig:entity-relationships}
\end{figure}

Фигура \ref{fig:entity-relationships} представя опростено отношенията между основните трайни обекти. От нея ясно се вижда, че потребителят е централният домейн елемент, към който се асоциират фирмите, количката и историята, докато продуктите са свързани с фирмите и се използват в количката.

\subsection{Атрибутен анализ на основните entity класове}

За по-систематичен преглед е полезно ключовите класове да бъдат сравнени по техните атрибути и роли в системата.

\begin{longtable}{L{2.7cm}L{3.8cm}L{7.3cm}}
\caption{Атрибути и роля на основните entity класове}\label{tab:entity-attributes}\\
\toprule
Клас & Основни атрибути & Роля в системата \\
\midrule
\endfirsthead
\toprule
Клас & Основни атрибути & Роля в системата \\
\midrule
\endhead
\texttt{User} & \texttt{firstName}, \texttt{lastName}, \texttt{password}, \texttt{email}, \texttt{isConfirmed} & Представя идентичността на потребителя и служи като входна точка към фирми, количка и история \\
\texttt{Company} & \texttt{name}, \texttt{mol}, \texttt{vat}, \texttt{address}, \texttt{isRegistered} & Представя фирмен субект, който се асоциира с потребителя и към който принадлежат продукти \\
\texttt{Product} & \texttt{name}, \texttt{price}, \texttt{quantity}, \texttt{url}, \texttt{category} & Представя каталожен артикул, който може да бъде визуализиран, филтриран и закупуван \\
\texttt{ShoppingCartEntity} & \texttt{wishedQuantityForOrder}, връзка към продукт и потребител & Представя временно потребителско намерение за покупка \\
\texttt{HistoryEntity} & \texttt{address}, \texttt{date}, \texttt{price}, връзка към потребител & Представя обобщение на финализирана поръчка \\
\bottomrule
\end{longtable}

Таблица \ref{tab:entity-attributes} потвърждава, че домейн моделът е достатъчно компактен, но и достатъчно изразителен. В него няма излишни слоеве на абстракция, а всеки entity клас има ясно разпознаваема роля.

\subsection{Изброим тип за категориите}

Категоризационният модел е реализиран чрез \texttt{CategoryEnum}. В него са дефинирани стойности като \texttt{VIDEOCARDS}, \texttt{PROCESSORS}, \texttt{HEADPHONES}, \texttt{PRINTERS}, \texttt{OFFICE\_CHAIRS}, \texttt{SCANNERS}, \texttt{ROUTERS}, \texttt{MICE} и \texttt{KEYBOARDS}. Този избор е особено подходящ по три причини.

Първо, елиминира се рискът от невалидни категории, въведени чрез свободен текст. Второ, контролерът за каталог може да използва ясни branch условия по категории. Трето, категориите се превръщат в част от самия домейн модел, а не остават външен конфигурационен списък без типова защита.

Цената на това решение е, че добавянето на нова категория изисква промяна в кода и повторна компилация. За текущия проект това не е проблем, тъй като категориите са предварително известни и тяхната стабилност е по-важна от динамичната конфигурируемост.

\subsection{Binding модели и входна валидация}

Съществена част от домейн обработката в системата не се извършва директно чрез entity класовете, а чрез binding модели. Това са класове, предназначени за приемане на данни от HTML форми. Сред тях са \texttt{UserRegisterBindingModel}, \texttt{UserLoginBindingModel}, \texttt{ProfileUpdateBindingModel}, \texttt{AddCompanyBindingModel}, \texttt{ProductAddBindingModel}, \texttt{AddToCartBindingModel} и \texttt{CheckoutBindingModel}.

От архитектурна гледна точка тези класове изпълняват две роли. Първо, дефинират структурата на потребителския вход. Второ, концентрират на едно място част от валидационната логика чрез анотации като \texttt{@Size}, \texttt{@NotBlank}, \texttt{@NotNull} и \texttt{@Positive} \cite{hibernatevalidator}. Това прави формите по-прогнозируеми и намалява вероятността до service слоя да достигнат очевидно невалидни данни.

\begin{table}[H]
    \centering
    \caption{Примери за валидационни ограничения в binding моделите}
    \label{tab:binding-validation}
    \begin{tabularx}{\textwidth}{L{4.1cm}L{4cm}Y}
        \toprule
        Модел & Ограничения & Смисъл на ограничението \\
        \midrule
        \texttt{ProfileUpdateBindingModel} & \texttt{@Size(min=3, max=20)} за име, фамилия, имейл и пароли & Ограничаване на очевидно невалидни или прекомерно кратки стойности \\
        \texttt{AddCompanyBindingModel} & \texttt{@NotBlank}, \texttt{@Size}, \texttt{@Positive}, \texttt{@NotNull} & Валидиране на фирмено име, адрес и числови идентификатори \\
        \texttt{ProductAddBindingModel} & \texttt{@NotBlank}, \texttt{@Positive}, \texttt{@NotNull} & Гарантира коректни данни за продукт преди запис \\
        \texttt{CheckoutBindingModel} & \texttt{@Size(min=5, max=20)} за адрес & Базова проверка на адреса при финализиране на поръчка \\
        \bottomrule
    \end{tabularx}
\end{table}

Макар този модел да е полезен, следва да се отбележи, че не всички binding класове в проекта имат еднакво богата валидация. Например регистрационният и login моделът са по-леки и не носят симетричен набор от ограничения. Това е напълно допустимо в учебна разработка, но би могло да бъде разширено при бъдещо развитие.

\subsection{Service и view модели}

Освен entity и binding класове проектът използва и service модели, и view модели. Service моделите служат като вътрешно представяне между persistence и service слоя, а view моделите са оптимизирани за визуализация в шаблоните или за връщане през API. Примери за такива класове са \texttt{UserServiceModel}, \texttt{UserLoginServiceModel}, \texttt{CompanyServiceModel}, \texttt{CompanyViewModel}, \texttt{ProductViewModel}, \texttt{UserViewModel} и \texttt{MaxPriceViewModel}.

Предимството на този подход е, че отделните слоеве могат да работят с различни представяния на едни и същи логически данни, без да бъдат принудени да използват директно entity обектите навсякъде. Това подобрява капсулацията. Например шаблонът за каталог не е длъжен да получава пълния entity обект на продукта, а само тази част от информацията, която е необходима за визуализация.

\subsection{Преобразувания между моделите}

Преобразуването между различните представяния е подпомогнато от \texttt{ModelMapper}. Тази библиотека намалява ръчния код при преобразуване на entity към service модел или на service модел към view модел. В контекста на проекта това улеснява реализацията на списъци с потребители, фирми и продукти. Въпреки това част от домейн логиката остава ръчно реализирана, например когато е необходимо да се създаде изцяло нов entity обект с попълване на конкретни полета и зависимости.

Използването на mapper инструмент е особено полезно при средни по размер проекти. То намалява повторяемия код, но същевременно не скрива домейн решенията, когато те са специфични. Именно такъв е случаят тук: стандартните трансформации се автоматизират, а смислово важните операции се извършват експлицитно в service слоя.

\subsection{Съответствие между домейн модел и потребителски сценарии}

Една от силните страни на разглеждания проект е, че домейн моделът сравнително пряко следва основните потребителски сценарии. Потребителят създава фирма; фирмата притежава продукти; продуктът се добавя в количката; количката участва във финализиране на поръчка; поръчката се архивира в историята. Тази линия е достатъчно ясна, за да бъде следвана както от кода, така и от документацията.

Разбира се, моделът не е напълно независим от ограниченията на имплементацията. Историята не съдържа отделни позиции по поръчка, а количката е опростена спрямо типична production система. Въпреки това домейнът е достатъчно консистентен, за да изпълни предназначението си като основа на учебно-приложен онлайн магазин.

\subsection{Критични наблюдения върху модела}

Анализът на модела позволява да бъдат формулирани и няколко критични наблюдения:

\begin{itemize}[leftmargin=*]
    \item налице са дребни несъответствия в именуването на някои свойства, което затруднява концептуалната чистота;
    \item поръчката е представена в силно опростен вид чрез исторически запис, без отделни order item обекти;
    \item някои валидационни правила са концентрирани само върху част от binding моделите;
    \item моделът е силно ориентиран към един потребителски тип и не включва роли или права.
\end{itemize}

Тези наблюдения не обезсилват решението, а по-скоро очертават границите на неговата текуща зрялост. Именно поради това те са ценни за документацията: показват, че системата е анализирана критично, а не само описателно.

\subsection{Обобщение}

Домейн моделът на разглежданото приложение е компактен, но структуриран достатъчно ясно, за да поддържа основните процеси на електронен магазин. Чрез комбинация от entity обекти, binding модели, service модели и view модели проектът реализира последователен поток на данните от интерфейса до базата данни и обратно. Този модел е една от ключовите причини системата да бъде методологично подходяща за дипломна работа: той съчетава достатъчна практическа реалистичност с достатъчна архитектурна четимост. Следващата глава разглежда как този модел оживява в реалната back-end имплементация на контролерите и услугите.
